\documentclass[12pt]{article}
\usepackage[utf8]{inputenc}
\usepackage[T1]{fontenc}
\usepackage{hyperref}
\title{Talieng Language Resources for K-12 Curriculum Development}
\author{Compiled by Claude Sonnet 5 (claude-sonnet-5) for Andr\'as Kornai}
\date{}
\begin{document}
\maketitle
\section{Introduction}
Talieng (ISO 639-3 tdf, Glottocode tali1257) is spoken in Laos. This report is based on one web search. The sources found classify it as Bahnaric, within Austroasiatic, and give the alternative names Taliang, Tariang and Trieng; the Wikipedia article ``Tariang language'' reports 45,000 native speakers (2005). I did not verify that these are all the same language. Very little was found.
\section{Primary Educational Resources}
None found.
\section{Linguistic and Reference Materials}
\begin{itemize}
\item Numeral systems page for Talieng, MPI-EVA: \url{https://lingweb.eva.mpg.de/channumerals/Talieng.htm}
\item Wikipedia, Tariang language: \url{https://en.wikipedia.org/wiki/Tariang_language}
\item Glottolog bookkeeping note on Tareng (tgr), which discusses a Katuic/Bahnaric mix-up: \url{https://glottolog.org/resource/languoid/iso/tgr}
\end{itemize}
\section{Religious and Community Resources}
None found.
\section{Digital and Audio Resources}
None found.
\section{Curriculum Development Context}
The search found no teaching materials, orthography documentation or community organisations. A reliable curriculum would require fieldwork or searches in Lao-language and Vietnamese sources.
\end{document}